\chapter{Conclusion}
\label{ch:conclusion}

\begin{itemize}
\item \textbf{Formulation and implementation.} The fleet design problem was finalised as a
mixed-integer, single-objective problem in which the optimiser chooses the number of spacecraft
(on/off switches), the circular orbit of each, and the mission duration, trading coverage of the
sunlit surface against fleet cost, energy and time under radius, separation and battery
constraints. It is implemented as a modular Algorithm $\rightarrow$ Feasibility $\rightarrow$ Engine
pipeline around a true-scale simulation (Kepler orbits, spinning non-convex shape model, occlusion,
sunlight and shadows, camera resolution, batteries). Both \ac{SA} and \ac{GA} use this same pipeline;
their search operators remain separate from the physics model.
\item \textbf{Designs.} \ac{SA} sends only the spacecraft a mission needs: two for long Eros missions,
which it ends after about two days at 95\,\% coverage, and one to five (about three on average) when
only hours are available. The same code and settings work unchanged on the 46 times larger, strongly concave
Kleopatra (two to three spacecraft, 75\,\% in about 18\,h). The optimiser discovered on its own that the
best orbits sit just above the camera's range, at 3.5 times the body's radius on both bodies, and are
tilted about 30--35$^\circ$ rather than polar because the poles are never lit.
\item \textbf{Validation.} The simulator agrees with hand calculations; \ac{SA} finds the known optima of
benchmark functions; it beats random search with the same budget in all nine settings and in every
run ($p=0.004$), by 6 to 21 points of coverage, the more so the shorter the mission; under the
assumptions of the literature it images 98\,\% of Eros, close to the 100\,\% reported there; its fleet
sizes lie in the flat optimum found by forcing the fleet size; and its designs hold on a four times
finer model.
\item \textbf{Algorithm.} The cooling schedule is the most important parameter: fitting $\alpha$ to the
budget avoids both premature freezing and a search that never settles. A moderate start temperature
($p_0\le0.5$) and large initial steps are best. The track cache makes a move cost at most two
spacecraft simulations, so \ac{SA} is also two to four times faster per evaluation than random search.
\item \textbf{Genetic Algorithm (Milestone 4).} The \ac{GA} uses the whole mission as chromosome,
aligns spacecraft slots before crossover, swaps or blends complete orbits, starts with a population
stratified over fleet size and reuses \ac{SA}'s local move as mutation. A 23-setting parameter study
confirmed the default design choices and showed that both operator families are needed. On the
nine case studies \ac{SA} is slightly better (mean $J$ \saMeanJ{} against \gaMeanJ, significant in
\saWins{} cases) and converges faster, while the \ac{GA} is on a par in \gaTies{} cases, varies less
between seeds, has the lower mean fitness on the shortest Eros and Kleopatra missions and profits
more from a larger budget
(significantly better on Eros 2\,h with 5000 evaluations). In the matched 51-variable Eros case,
\ac{SA} achieved mean $J=\saMatchedJ$ and \saMatchedC\,\% coverage, compared with $J=\gaMatchedJ$ and
\gaMatchedC\,\% for the \ac{GA}; \ac{SA} was faster and used fewer spacecraft on average, while the
\ac{GA} produced a shorter mean mission. A hybrid local-refinement stage is the clearest next improvement for the \ac{GA}.
\item \textbf{Limitations.} The orbits ignore the irregular gravity field; the Sun direction is fixed;
the weights decide the value of one more spacecraft or day; and the current comparisons remain based
on a finite number of random seeds
(Chapter~\ref{ch:future}).
\end{itemize}
